\documentclass[11pt,letterpaper]{article}
\usepackage[T1]{fontenc}
\usepackage[utf8]{inputenc}
\usepackage[margin=0.88in,top=0.84in,bottom=0.83in,headheight=13pt,headsep=20pt,footskip=28pt]{geometry}
\usepackage{newpxtext}
\usepackage{amsmath,amsthm}
\usepackage{newpxmath}
\usepackage{microtype}
\usepackage{xcolor}
\definecolor{Forest}{HTML}{30392C}
\definecolor{Olive}{HTML}{687144}
\definecolor{Muted}{HTML}{65685F}
\definecolor{Sage}{HTML}{CBD0BC}
\definecolor{Pale}{HTML}{F2F4EB}
\usepackage{mathtools}
\usepackage{booktabs,tabularx,longtable,array}
\usepackage{enumitem}
\usepackage{titlesec}
\usepackage{fancyhdr}
\usepackage[most]{tcolorbox}
\usepackage{needspace}
\PassOptionsToPackage{obeyspaces}{url}
\usepackage{xurl}
\usepackage[colorlinks=true,linkcolor=Olive,citecolor=Olive,urlcolor=Olive,bookmarksnumbered=true,pdfencoding=auto,psdextra]{hyperref}
\hypersetup{pdftitle={Finite choice at measurable strength: Prikry deletion and profinite phases},pdfauthor={Prepared for Vladimir Reshetnikov},pdfsubject={A continuation of Cardinals4: finite-family cyclic obstructions in an ordinary Prikry extension},pdfkeywords={Prikry forcing, measurable cardinal, finite choice, HOD, cyclic stabilizer, profinite integers}}
\urlstyle{same}
\setlength{\parindent}{0pt}
\setlength{\parskip}{5.1pt plus 1pt minus .4pt}
\linespread{1.035}
\setlength{\emergencystretch}{2em}
\setcounter{tocdepth}{1}
\setcounter{secnumdepth}{2}
\titleformat{\section}{\Large\bfseries\color{Forest}}{\thesection}{.6em}{}
\titleformat{\subsection}{\large\bfseries\color{Forest}}{\thesubsection}{.55em}{}
\titleformat{\subsubsection}{\normalsize\bfseries\color{Forest}}{}{0pt}{}
\titlespacing*{\section}{0pt}{18pt plus 3pt minus 2pt}{8pt}
\titlespacing*{\subsection}{0pt}{12pt plus 2pt minus 1pt}{5pt}
\titlespacing*{\subsubsection}{0pt}{9pt}{3pt}
\setlist[itemize]{leftmargin=1.4em,itemsep=2pt,topsep=3pt,parsep=0pt}
\setlist[enumerate]{leftmargin=1.7em,itemsep=3pt,topsep=4pt,parsep=0pt}
\pagestyle{fancy}
\fancyhf{}
\fancyhead[L]{\sffamily\fontsize{9}{11}\selectfont\color{Muted}LARGE CARDINALS AT THE INCONSISTENCY FRONTIER}
\fancyhead[R]{\sffamily\fontsize{9}{11}\selectfont\color{Muted}CONTINUATION 4}
\fancyfoot[L]{\small\color{Muted}Prikry deletion and finite choice}
\fancyfoot[R]{\small\color{Muted}\thepage}
\renewcommand{\headrulewidth}{.35pt}
\renewcommand{\footrulewidth}{0pt}
\fancypagestyle{plain}{\fancyhf{}\fancyfoot[L]{\small\color{Muted}Prikry deletion and finite choice}\fancyfoot[R]{\small\color{Muted}\thepage}\renewcommand{\headrulewidth}{0pt}}
\tcbset{colback=Pale,colframe=Sage,colbacktitle=Sage,coltitle=Forest,fonttitle=\bfseries,boxrule=.5pt,arc=3pt,left=9pt,right=9pt,top=6pt,bottom=6pt,toptitle=3pt,bottomtitle=3pt,before skip=8pt,after skip=8pt,breakable}
\newtcolorbox{assessment}[1]{title={#1}}
\theoremstyle{definition}
\newtheorem{definition}{Definition}[section]
\newtheorem{theorem}[definition]{Theorem}
\newtheorem{proposition}[definition]{Proposition}
\newtheorem{lemma}[definition]{Lemma}
\newtheorem{corollary}[definition]{Corollary}
\newtheorem{remark}[definition]{Remark}
\newtheorem{question}[definition]{Question}
\newtheorem{inputthm}{Input}
\renewcommand{\theinputthm}{\Alph{inputthm}}
\numberwithin{equation}{section}
\newcommand{\ZFC}{\mathsf{ZFC}}
\newcommand{\ZF}{\mathsf{ZF}}
\newcommand{\AC}{\mathsf{AC}}
\newcommand{\GA}{\mathsf{GA}}
\newcommand{\HOD}{\mathrm{HOD}}
\newcommand{\OD}{\mathrm{OD}}
\newcommand{\CD}{\mathrm{CD}}
\newcommand{\HCD}{\mathrm{HCD}}
\newcommand{\UF}{\mathrm{UF}}
\newcommand{\Ex}{\operatorname{Ex}}
\newcommand{\UEx}{\operatorname{UEx}}
\newcommand{\Ext}{\operatorname{Ext}}
\newcommand{\SC}{\operatorname{SC}}
\newcommand{\CEx}{\operatorname{CEx}}
\newcommand{\REx}{\operatorname{REx}}
\newcommand{\Con}{\operatorname{Con}}
\newcommand{\crit}{\operatorname{crit}}
\newcommand{\cf}{\operatorname{cf}}
\newcommand{\ot}{\operatorname{ot}}
\newcommand{\ran}{\operatorname{ran}}
\newcommand{\rank}{\operatorname{rank}}
\newcommand{\lcm}{\operatorname{lcm}}
\newcommand{\Ord}{\mathrm{Ord}}
\newcommand{\Pow}{\mathcal P}
\newcommand{\id}{\mathrm{id}}
\newcommand{\restr}{\mathbin{\upharpoonright}}
\newcommand{\Izero}{\mathrm{I0}}
\newcommand{\Itwo}{\mathrm{I2}}
\newcommand{\Ithree}{\mathrm{I3}}
\newcommand{\Iwf}{\mathrm{I3}_{\mathrm{wf}}(0)}
\newcommand{\Sel}{\operatorname{Sel}}
\newcommand{\FF}{\mathcal F}
\newcommand{\PP}{\mathbb P}
\newcommand{\Clam}{\mathcal C_\lambda}
\newcommand{\Rig}{\operatorname{Rig}}
\newcommand{\Spec}{\operatorname{Spec}}
\newcommand{\ch}{\operatorname{ind}}
\newcommand{\CF}{\mathsf{CF}}
\newcommand{\src}[1]{\unskip\ {\normalfont\small\sffamily\color{Olive}[#1]}\ }
\newcommand{\R}[1]{\textsf{R#1}}
\newcolumntype{Y}{>{\raggedright\arraybackslash}X}
\newcolumntype{P}[1]{>{\raggedright\arraybackslash}p{#1}}
\newcolumntype{C}{>{\centering\arraybackslash}p{.034\textwidth}}
\newcommand{\yes}{$\bullet$}
\newcommand{\prt}{$\circ$}
\newcommand{\arxiv}[1]{\href{https://arxiv.org/abs/#1}{arXiv:#1}}
\newcommand{\doi}[1]{\href{https://doi.org/#1}{doi:\nolinkurl{#1}}}


\DeclareMathOperator{\Sym}{Sym}
\DeclareMathOperator{\Inj}{Inj}
\DeclareMathOperator{\Fix}{Fix}
\DeclareMathOperator{\Stab}{Stab}
\DeclareMathOperator{\Ult}{Ult}
\newcommand{\ZZ}{\mathbb Z}
\newcommand{\NN}{\mathbb N}
\newcommand{\RR}{\mathbb R}
\newcommand{\zhat}{\widehat{\mathbb Z}}
\newcommand{\Dk}{D_\kappa}
\newcommand{\Qk}{Q_\kappa}
\newcommand{\ODg}{\OD_{M\cup\{q\}}^{V}}
\newcommand{\ODlow}{\OD_{V_\kappa\cup\{q\}}^{V}}
\newcommand{\tail}{\operatorname{Tail}}
\newcommand{\Inv}{\mathsf{FC}}
\newcommand{\NTFC}{\mathsf{NT+FC}}
\newcommand{\powerset}{\mathcal P}
\newcommand{\cseq}{\langle c_n:n<\omega\rangle}
\begin{document}
\thispagestyle{empty}
{\sffamily\bfseries\small\color{Olive}SET THEORY\quad /\quad RESEARCH CONTINUATION}

\vspace{13pt}
{\fontsize{28}{31}\selectfont\bfseries\color{Forest}Finite choice at\ measurable strength\par}
\vspace{10pt}
{\fontsize{14}{18}\selectfont Prikry deletion, cyclic label obstructions,\ and profinite phases\par}
\vspace{14pt}
{\color{Muted}Prepared for Vladimir Reshetnikov\hfill 18 September 2026\par}
\vspace{3pt}
{\color{Sage}\hrule height .6pt}
\vspace{12pt}

\textbf{Abstract.}
Let $M\models\ZFC$, let $U\in M$ be a normal measure on $\kappa$, and let
$V=M[c]$ be its ordinary Prikry extension. Put
$q=[\ran(c)]$ modulo finite symmetric difference. We prove that no nonempty
finite family of selectors on the cofinal $\omega$-quotient $Q_\kappa$ is
ordinal definable in $V$ from $q$ and ground-model parameters. More generally,
for a finite permutation group $\Gamma\le S_m$ acting on a nonempty ground-model
set $Y$, a definable admissible labeling rule exists exactly when every
$g\in\Gamma$ fixes some element of $Y$. Allowing a nonempty finite definable
family of rules does not weaken this criterion.

The necessity proof uses a single-point deletion of the generic sequence,
a finite permutation root, and a finite table of responses. Thus it does not
require an ultraexacting cardinal or an elementary self-embedding. A second
result shows that every nonempty definable family of ultrafilters on the
\emph{set of representatives} $q$ has a phase image invariant under translation
by $1$ in $\zhat$, hence dense. A compact such family has at least $2^{2^{\aleph_0}}$
members, and this bound is attained. An explicit countably infinite definable family realizes
a single dense integer orbit, making the local finite obstruction sharp.
We also give sharp finite-additivity comparisons and combine
the new obstructions with the existing normal-trace construction to obtain
an exact measurable-strength consistency package.

\begin{assessment}{Main advance relative to the supplied archive}
The Prikry-extension part of the last open-question group in the supplied
\emph{Synthesis} has an affirmative answer: the finite-selector obstruction
and the necessity of its cyclic label criterion already hold after Prikry
forcing at one measurable cardinal. The finite-family label theorem and the
profinite-density conclusion strengthen that answer. A countable
ultrafilter construction also gives the exact least size, $\aleph_0$, of a
nonempty definable family of ultrafilters on the generic class.

This is \emph{not} a proof that rigidity alone implies these properties in
an arbitrary universe. Nor does it establish an inconsistency of
ultraexactingness, $\Izero$, or any other large-cardinal axiom.
\end{assessment}

\textbf{Proof status.}
The report supplies conventional proofs in English. Classical Prikry forcing
facts are identified as inputs; the new deductions are proved below.
The supplied Lean files were used as reference, not as a verification of
these forcing arguments. The additional Python checks concern only finite
permutation identities. Novelty is assessed relative to the supplied
reports; an exhaustive publication-priority claim is not made.

\newpage
\tableofcontents

\section{What is new, and what is being continued}\label{sec:scope}

The supplied archive contains the research plan, a synthesis of twenty-seven
continuations, and Lean reference material. The immediate source for this
report is \cite{archive}, particularly its sections on finite selectors,
cyclic stabilizers, internal normal measures, and exact consistency strength.
The final open-question group asks whether the finite-label necessity theorem
and the no-finite-family-of-selectors theorem hold in a Prikry extension of a
measurable cardinal. Those are the questions answered here.

The same sentence in the synthesis also asks whether the conclusions follow
from rigidity alone. The two formulations should be separated. The assertion
that a particular forcing construction has a property is weaker than the
assertion that every instance of an abstract rigidity principle has it. Our
proof uses a feature of a generic presentation that is not part of the
definition of rigidity: a sufficiently late deletion preserves both the
genericity and the truth of a previously decided ground-valued observation.

\begin{center}
\small
\renewcommand{\arraystretch}{1.18}
\begin{tabularx}{\textwidth}{@{}P{.29\textwidth}YY@{}}
\toprule
\textbf{Topic}&\textbf{Already in the synthesis}&\textbf{Contribution here}\\
\midrule
Finite selectors&No finite definable family at an ultraexacting cardinal&The same obstruction in ordinary Prikry extensions, even with the generic class and arbitrary ground parameters\\
Cyclic label criterion&Single-rule classification for finite label sets&Finite-family necessity in Prikry extensions; ground label sets may be infinite\\
Ultrafilter-valued choice&Finite-family obstruction at ultraexacting cardinals&Dense phase images in $\zhat$; exact local thresholds $\aleph_0$ and, under compactness, $2^{2^{\aleph_0}}$\\
Rigidity and tail measures&Measurable-strength calibration and Prikry cores&A full compatibility proof incorporating all finite cyclic obstructions at the same strength\\
Positive label rules&Incidence words with cyclic stabilizers&Reproduced with details, not claimed as a new method\\
\bottomrule
\end{tabularx}
\end{center}

The broader motivation remains the analysis of exacting and ultraexacting
cardinals in \cite{abl,abgl}. Their embedding arguments explain why these
quotient obstructions arose in the archive. The present proof isolates a
forcing mechanism sufficient for them. The use of finite cycles in quotient
choice arguments also has precedent in the discussion recorded in
\cite[Remark 2.4]{goldberg}; the result here is the Prikry transfer and its
finite-family amplification, not the observation that a cycle can obstruct
a choice of one of its elements.

\begin{assessment}{Reading guide}
The main new proof is in Sections~\ref{sec:deletion}--\ref{sec:necessity}.
Section~\ref{sec:sufficiency} proves the converse and gives the complete
classification. Sections~\ref{sec:phases}--\ref{sec:charges} contain the
ultrafilter and probability refinements. Sections~\ref{sec:core} and
\ref{sec:consistency} give the inner-model comparison and the precise
consistency statement.
\end{assessment}

\section{The setting and the classical forcing inputs}\label{sec:setting}

\subsection{Models, quotients, and parameters}

Throughout Sections~\ref{sec:setting}--\ref{sec:core}, $M$ is a transitive
model of $\ZFC$, $U\in M$ is a normal nonprincipal $\kappa$-complete
ultrafilter on $\kappa$ as computed in $M$, and $c=\cseq$ is increasing and
Prikry generic for $U$ over $M$. We write
\[
 V=M[c],\qquad a_c=\ran(c),\qquad q=[a_c].
\]
The same arguments may be read as forcing arguments over a ground universe.
The metamathematical consistency statement at the end does not assume a
transitive model exists.

For an uncountable ordinal $\lambda$ with $\cf(\lambda)=\omega$, define
\[
 D_\lambda=\{a\subseteq\lambda:\sup a=\lambda\text{ and }\ot(a)=\omega\},
 \qquad
 a=^*b\ \Longleftrightarrow\ |a\mathbin\triangle b|<\omega,
 \qquad Q_\lambda=D_\lambda/=^*.
\]
All these sets, and all unqualified cardinalities below, are computed in
$V$. A representative and its increasing enumeration will be distinguished
where indices matter. For $q\in Q_\lambda$ and $A\subseteq\lambda$, put
\[
 \tail(q,A)\ \Longleftrightarrow\ a\setminus A\text{ is finite for some }a\in q.
\]
``Some'' can equivalently be replaced by ``every.''

\begin{definition}[Parameter convention]\label{def:parameters}
The notation $\ODg$ means the sets uniquely definable in $V$ from $q$,
finitely many parameters belonging to $M$, and finitely many ordinals.
It does \emph{not} allow the class $M$ as a predicate in the definition.
Similarly, $\ODlow$ allows finitely many elements of $V_\kappa$ and the
single parameter $q$. In particular, $q$ is allowed but its representatives
are not automatically allowed.
\end{definition}

A finite family $\mathcal F$ being definable does not say that a member of
$\mathcal F$ is definable. Our main theorem is deliberately stated for
families and does not make that invalid reduction.

\subsection{The forcing and its standard properties}

A condition of $\PP_U$ is a pair $(s,A)$, where $s$ is a finite increasing
sequence in $\kappa$, $A\in U$, and every element of $A$ is above the stem
$s$. Stronger conditions extend the stem by points from $A$ and shrink
$A$. We use the convention that a smaller condition is stronger.

\begin{inputthm}[Classical Prikry facts]\label{input:prikry}
Normal Prikry forcing preserves cardinals, adds no bounded subset of
$\kappa$, and changes $\cf(\kappa)$ to $\omega$. An increasing cofinal
$\omega$-sequence $d$ in $\kappa$ is generic over $M$ precisely when it is
almost contained in every member of $U$. For such $d$, its generic filter is
\[
 G_d=\{(s,A):s\text{ is an initial segment of }d
                  \text{ and }\ran(d)\setminus\ran(s)\subseteq A\}.
\]
\end{inputthm}

The characterization is Mathias' criterion. The input is classical; see
\cite{mathias} and the modern proofs in \cite[Section 3, especially
Corollary 3.15 and Theorem 3.17]{benhamou}. We use the forcing theorem and
truth lemma in their usual set-forcing form. Everything below the stated
input is proved here, except where an earlier archive construction is
explicitly credited.

\begin{lemma}[No new low-rank sets]\label{lem:lowrank}
$V_\kappa^V=V_\kappa^M$. Consequently $\ODlow\subseteq\ODg$; the two
models have the same reals and the same sets of reals. Moreover, $\kappa$
is an uncountable strong limit cardinal of cofinality $\omega$ in $V$.
\end{lemma}

\begin{proof}
The cardinal $\kappa$ is strongly inaccessible in $M$. We prove equality
of $V_\alpha$ for $\alpha<\kappa$ by induction. The limit step is immediate.
At a successor step, a ground bijection codes $V_\alpha^M$ by an ordinal
below $\kappa$, since $|V_\alpha^M|^M<\kappa$. A new subset of this set
would therefore code a new bounded subset of $\kappa$, contrary to
Input~\ref{input:prikry}. This proves the induction and the equality.
Reals, sets of reals, and their standard finite-rank codes all lie below
$V_\kappa$. Finally cardinals are preserved, and every power set of an
ordinal below $\kappa$ is unchanged and has cardinality below $\kappa$.
Thus $\kappa$ remains strong limit, while its cofinality becomes $\omega$.
\end{proof}

\begin{lemma}[Finite changes preserve the presentation]\label{lem:finitechange}
If $d$ is the increasing enumeration of any member of $q$, then $d$ is
Prikry generic over $M$ and $M[d]=M[c]$. In particular, deleting $c_k$ from
$c$ gives such a $d$. If $(s,A)\in G_c$ and $k\ge |s|$, then
$(s,A)\in G_d$ for this deletion.
\end{lemma}

\begin{proof}
A finite change preserves cofinality and eventual membership in each
$A\in U$, so genericity follows from Mathias' criterion. The two sequences
are interdefinable from their finite symmetric difference and finitely many
ordinals. All those additional parameters belong to $M$. Hence the two
extensions contain one another. For the last assertion the deletion does
not change the initial stem $s$ and does not add a point outside $A$;
the displayed description of $G_d$ applies.
\end{proof}

\section{The deletion principle for ground-valued observations}\label{sec:deletion}

The following lemma is the central forcing tool. It does not refer to an
elementary embedding and does not assert that $V$ has a nontrivial
automorphism.

\begin{lemma}[A late deletion preserves a decided observation]\label{lem:observation}
Let $A\in M$. Suppose $\Phi$ is a fixed formula that, from an increasing
generic sequence $e$, its class $[\ran(e)]$, ground parameters, and ordinal
parameters, uniquely specifies an element $h(e)\in A$ in the generic
extension. It is enough to assume this uniqueness in the presentations
under consideration, including $c$ and its finite changes. Then there is
$k_0<\omega$ such that for every $k\ge k_0$, if $d$ is obtained from $c$
by deleting its $k$th term, then
\[
 h(d)=h(c).
\]
In particular, if a fixed ground function $T:A\to A$ satisfies
$h(d)=T(h(c))$ for every sufficiently late deletion, then
$T(h(c))=h(c)$.
\end{lemma}

\begin{proof}
Write $a=h(c)$. Since $a\in A$ and $A\in M$ with $M$ transitive, $a\in M$.
By the truth lemma there is a condition $(s,B)\in G_c$ forcing that the
unique observation, formed using the canonical generic sequence and its
class, is the ground object $\check a$. This statement includes whatever
uniqueness assertion is needed for $\Phi$.

Take $k_0=|s|$. If $k\ge k_0$ and $d$ deletes $c_k$,
Lemma~\ref{lem:finitechange} shows that $G_d$ contains $(s,B)$ and that
$M[d]=V$. Therefore the very same forced definition, interpreted using
$d$, has value $a$. This proves $h(d)=a=h(c)$. The concluding fixed-point
identity follows by comparing the two formulas for $h(d)$.
\end{proof}

\begin{remark}[What the lemma does and does not fix]\label{rem:observation}
The object $h(e)$ is allowed to depend on the actual enumeration $e$;
otherwise the lemma would have little content. The definable family from
which the observation is extracted will depend only on ground parameters
and the unchanged class $q$, so that family is the same set in $M[c]$ and
$M[d]$. The observation is low-complexity enough to be a ground object,
although its inputs and the members of the family may be new and have high
rank. We compare two generic presentations of one extension, not images
under an automorphism of that extension.
\end{remark}

A useful example is a finite response table with entries in a ground set
$Y$. Every such table is in $M$, because $M$ is closed under forming finite
sets and ordered tuples of its elements. A finite set of these tables is
also in $M$. No definable enumeration of that finite set is required.
A second example, used in Section~\ref{sec:phases}, is an arbitrary subset
of a fixed real-coded compact space: this is a ground object by
Lemma~\ref{lem:lowrank}.

\section{Encoding an arbitrary finite permutation}\label{sec:coding}

\begin{lemma}[One deletion realizes one permutation]\label{lem:coding}
Let $E$ be a nonempty finite set, let $\tau\in\Sym(E)$, and fix an injective
tag map $t:E\to N$ for some positive integer $N$. There are representatives
$a_e(c)\in D_\kappa$, for $e\in E$, whose classes are pairwise distinct,
such that for every sufficiently late deletion $d$ of $c$,
\begin{equation}\label{eq:permutation}
 a_e(d)=^*a_{\tau(e)}(c)\qquad(e\in E).
\end{equation}
They can be defined uniformly from $c,\tau,t$ by
\begin{equation}\label{eq:coding}
 a_e(c)=\{c_n+t(\tau^{-n}(e)):n<\omega\}.
\end{equation}
\end{lemma}

\begin{proof}
Normality implies that $U$ contains every ground club in $\kappa$. In
particular $c_n$ is eventually a nonzero limit ordinal. For completeness,
the club assertion follows because on a measure-one subset of the
complement of a club $C$, the regressive map
$\alpha\mapsto\sup(C\cap\alpha)$ would have to be constant. A constant
value would bound that subset below the next point of $C$, contradicting
nonprincipality and $\kappa$-completeness.

For all sufficiently large $n$, the finite block
$\{c_n+i:i<N\}$ is below $c_{n+1}$, since $c_{n+1}$ is a limit ordinal.
The blocks are therefore eventually disjoint and in increasing order.
Formula~\eqref{eq:coding} chooses one point in each block. It defines a
cofinal set of order type $\omega$; any exceptional early blocks contribute
only finitely many points. For $e\ne f$, injectivity of $t$ and of
$\tau^{-n}$ gives distinct points in every sufficiently late block.
Thus $a_e(c)\cap a_f(c)$ is finite, so their classes are distinct.

Suppose $d$ deletes $c_k$. For $n\ge k$, $d_n=c_{n+1}$, and
\[
 d_n+t(\tau^{-n}(e))
 =c_{n+1}+t\bigl(\tau^{-(n+1)}(\tau(e))\bigr).
\]
This is the $(n+1)$st term of the defining sequence for $a_{\tau(e)}(c)$.
The initial discrepancies are finite, proving~\eqref{eq:permutation}.
\end{proof}

\begin{lemma}[A finite permutation root]\label{lem:root}
Let $g\in S_m$ and $L\ge1$. On $E=m\times L$ define
\begin{equation}\label{eq:root}
 \tau(i,a)=
 \begin{cases}
  (i,a+1),&a+1<L,\\
  (g(i),0),&a=L-1.
 \end{cases}
\end{equation}
Then $\tau$ is a permutation and
$\tau^L(i,0)=(g(i),0)$ for each $i<m$.
\end{lemma}

\begin{proof}
The inverse sends $(i,a)$ with $a>0$ to $(i,a-1)$ and sends $(i,0)$ to
$(g^{-1}(i),L-1)$. Starting at $(i,0)$, the first $L-1$ steps visit the
remaining levels over $i$, and the $L$th step reaches $(g(i),0)$.
This includes $L=1$.
\end{proof}

The root is the amplification device. A permutation of a family of size
$h$ need not fix each member, but its $L$th power does when
$L=\lcm(1,\ldots,h)$. We arrange that this $L$th power still acts as the
desired permutation $g$ on a distinguished input tuple.

\section{The finite-family cyclic obstruction}\label{sec:necessity}

Fix $1\le m<\omega$; the case $m=0$ has a trivial group and requires no argument.
Let $D_m(\kappa)$ consist of ordered $m$-tuples
$\vec a=(a_i)_{i<m}$ from $D_\kappa$ whose finite-difference classes are
pairwise distinct. If $\Gamma\le S_m$, its action is
\[
 (g\cdot\vec a)_i=a_{g^{-1}(i)}.
\]
Let $Y\in M$ be nonempty and let a left action of $\Gamma$ on $Y$ belong to
$M$. The set $Y$ need not be finite.

\begin{definition}[Admissibility and the cyclic condition]\label{def:admissible}
A rule $f:D_m(\kappa)\to Y$ is \emph{admissible} if it is invariant under
finite changes in each coordinate and satisfies
$f(g\cdot\vec a)=g\cdot f(\vec a)$ for every $g\in\Gamma$.
Write
\[
 \CF(\Gamma,Y)\quad\Longleftrightarrow\quad
       (\forall g\in\Gamma)(\exists y\in Y)\ g\cdot y=y.
\]
Equivalently, every cyclic subgroup of $\Gamma$ has a fixed point in $Y$.
The label $y$ may depend on the cyclic subgroup.
\end{definition}

\begin{theorem}[Finite-family necessity in a Prikry extension]\label{thm:necessity}
If a nonempty finite family $\mathcal F$ of admissible rules belongs to
$\ODg$, then $\CF(\Gamma,Y)$ holds.
\end{theorem}

\begin{proof}
Let $h=|\mathcal F|\ge1$ and $L=\lcm(1,\ldots,h)$. Fix $g\in\Gamma$.
We will obtain a fixed point of $g$ on $Y$.

\textbf{Step 1: enlarge the finite test configuration.}
Use Lemma~\ref{lem:root} to define $E=m\times L$ and $\tau\in\Sym(E)$.
Choose distinct natural-number tags for $E$, and use
Lemma~\ref{lem:coding} to obtain pairwise inequivalent representatives
$a_e(c)$ for $e\in E$. Let
\[
 I=\Inj(m,E),\qquad u_0(i)=(i,0).
\]
The set $I$ is finite and belongs to $M$.

\textbf{Step 2: replace each possibly high-rank rule by a finite table.}
For $f\in\mathcal F$, define its response table
\[
 T_f^c:I\longrightarrow Y,\qquad
 T_f^c(u)=f\bigl((a_{u(i)}(c))_{i<m}\bigr),
\]
and put
\[
 H(c)=\{T_f^c:f\in\mathcal F\}.
\]
The tables may identify different rules, but
$1\le |H(c)|\le h$. Every table is a finite tuple of ground objects,
so $H(c)\in M$. More uniformly, it belongs to the ground set
$[Y^I]^{\le h}$. The observation $H(c)$ is defined using $c$ and a fixed
definition of $\mathcal F$ from $q$, ground parameters, and ordinals.

\textbf{Step 3: one deletion permutes the table set.}
On $Y^I$, define the ground permutation
\[
 (P_\tau T)(u)=T(\tau\circ u).
\]
Delete a sufficiently late term from $c$ to obtain $d$. The family
$\mathcal F$ is unchanged, because $M[d]=M[c]$, the ground parameters are
unchanged, and $[\ran(d)]=q$. By~\eqref{eq:permutation} and
finite-modification invariance,
\[
 T_f^d(u)=T_f^c(\tau\circ u).
\]
Consequently $H(d)=P_\tau[H(c)]$. Lemma~\ref{lem:observation}, applied
to this ground-valued observation, gives
\begin{equation}\label{eq:Hfixed}
 H(c)=P_\tau[H(c)].
\end{equation}
This equality is an equality of \emph{sets of tables}; it does not presuppose
that a particular rule or table has been definably selected.

\textbf{Step 4: take a power that fixes every table.}
Equation~\eqref{eq:Hfixed} says that $P_\tau$ induces a permutation of the
nonempty finite set $H(c)$. Each of its cycles has length at most $h$,
hence its length divides $L$. Therefore, for every $T\in H(c)$,
\[
 P_\tau^L T=T,\qquad
 T(\tau^L\circ u)=T(u)\quad(u\in I).
\]
Choose any $T=T_f^c\in H(c)$. This is an ordinary existential choice in
the proof, not a definable selection. Applying the last identity at $u_0$
and using $\tau^L\circ u_0=u_0\circ g$ yields
\[
 T(u_0)=T(u_0\circ g).
\]
The tuple with coordinates $a_{u_0(g(i))}(c)$ is
$g^{-1}\cdot(a_{u_0(i)}(c))_{i<m}$. Equivariance therefore gives
\[
 T(u_0\circ g)=g^{-1}\cdot T(u_0).
\]
Thus $T(u_0)$ is fixed by $g$. Since $g\in\Gamma$ was arbitrary,
$\CF(\Gamma,Y)$ follows.
\end{proof}

\begin{remark}[Why two tempting shortcuts fail]
One cannot simply select the least member of $\mathcal F$: its members
need not be ordinal definable, and no suitable definable well-order of the
new rules has been given. Nor is it enough to repeat a singleton argument
for one chosen member: deletion may move that member's table to another
one. The finite response set and the $L$th-root construction are precisely
what eliminate both gaps.
\end{remark}

\subsection{Selectors and related structures}

For $1\le r<n<\omega$, a selector of type $(n,r)$ on $Q_\kappa$ is a function
\[
 s:[Q_\kappa]^n\longrightarrow[Q_\kappa]^r,
 \qquad s(B)\subseteq B.
\]

\begin{corollary}[No finite definable family of selectors]\label{cor:selectors}
For every $1\le r<n<\omega$, there is no nonempty finite family of
$(n,r)$-selectors on $Q_\kappa$ in $\ODg$. In particular this holds in
$\ODlow$ and in $\OD_{V_\kappa}^V$.
\end{corollary}

\begin{proof}
To a selector $s$ associate
\[
 f_s(\vec a)=\{i<n:[a_i]\in s(\{[a_j]:j<n\})\}\in[n]^r.
\]
This is invariant under finite modifications and is equivariant for the
natural $S_n$-action on $[n]^r$. A nonempty finite definable family of
selectors would give a nonempty finite definable family of these rules.
But an $n$-cycle acts transitively on $n$, so its only invariant subsets are
$\varnothing$ and $n$. It fixes no member of $[n]^r$ when $0<r<n$.
Theorem~\ref{thm:necessity} gives a contradiction. The parameter
inclusions follow from Lemma~\ref{lem:lowrank}.
\end{proof}

\begin{corollary}\label{cor:orders}
There is no nonempty finite $\ODg$ family of linear orders on $Q_\kappa$,
and none of tournaments on $Q_\kappa$. There is no $\ODg$ injection from
$Q_\kappa$ into the ordinals.
\end{corollary}

\begin{proof}
A linear order chooses the smaller of two distinct elements. A tournament
chooses the source of the unique directed edge between them. These
conversions send a finite definable family to a nonempty finite definable
family of binary selectors. An injection into the ordinals similarly
chooses the element with smaller image. Apply
Corollary~\ref{cor:selectors} with $(n,r)=(2,1)$.
\end{proof}

\begin{assessment}{The obstruction is definability, not Choice}
The ambient extension $V$ satisfies $\ZFC$. It has selectors and well-orders
of $Q_\kappa$. The conclusion is that these cannot be captured by a
nonempty finite family defined from the specified ground parameters and
$q$. Adding an individual generic representative as a parameter would
invalidate the deletion argument; it is not an allowed substitution for
$q$.
\end{assessment}

\section{The converse: cyclic stabilizers of incidence words}\label{sec:sufficiency}

This section expands the positive construction already present in the
synthesis \cite[Section 12.6]{archive}. The finite-family necessity theorem
above is new relative to that source; the incidence-word method is not.
We include the proof to make the resulting classification self-contained.

\subsection{Shift-tail equivalence}

Let $\vec a\in D_m(\kappa)$. The union $\bigcup_{i<m}a_i$ has order type
$\omega$: below any ordinal smaller than $\kappa$, it has only finitely
many points. Enumerate it increasingly as $\langle\delta_n:n<\omega\rangle$
and form the incidence word
\[
 w_{\vec a}(n)=\{i<m:\delta_n\in a_i\}
       \in\powerset(m)\setminus\{\varnothing\}.
\]
Each label $i$ appears infinitely often. Each pair of distinct labels is
separated infinitely often, because otherwise the corresponding
representatives would have finite symmetric difference.

For two such words, write $w\sim v$ when there are $r,s<\omega$ such that
\[
 w(r+n)=v(s+n)\qquad(n<\omega).
\]
This is an equivalence relation. Changing finitely many points of the
coordinates changes the incidence word only within its $\sim$-class:
beyond the largest changed ordinal the two unions, and their incidence
labels, agree, possibly with different starting indices. Also
$w_{g\cdot\vec a}=g\cdot w_{\vec a}$, where $g$ acts on subsets of $m$
by image.

\begin{lemma}[The stabilizer is cyclic]\label{lem:stabilizer}
Let $w$ be an incidence word in which any two labels are separated
infinitely often. Then the stabilizer of $[w]_\sim$ in $\Gamma$ is cyclic.
\end{lemma}

\begin{proof}
Define
\[
 G_w=\{(d,g)\in\ZZ\times\Gamma:
        w(n+d)=g\cdot w(n)\text{ for all sufficiently large }n\}.
\]
Negative values of $d$ cause no problem because $n$ is taken sufficiently
large. The identity belongs to $G_w$. Composing the eventual identities
for $(d,g)$ and $(e,h)$ gives the identity for $(d+e,gh)$, and reversing
one gives the inverse identity. Thus $G_w$ is a subgroup.

The projection $G_w\to\ZZ$ is injective. Indeed, if $(0,g)\in G_w$ and
$g$ moves $i$, then for all sufficiently large $n$,
\[
 i\in w(n)\quad\Longleftrightarrow\quad g^{-1}(i)\in w(n),
\]
contradicting infinite separation of those two labels. Hence the kernel
is trivial. A subgroup of $\ZZ$ is cyclic: a nonzero subgroup is generated
by its least positive element, and the zero subgroup is trivial.
Therefore $G_w$ is cyclic. Its image in $\Gamma$ is exactly the stabilizer
of $[w]_\sim$, since an equality of two shifted tails is precisely an
eventual equality with an integer shift. An image of a cyclic group is
cyclic, proving the claim.
\end{proof}

\begin{theorem}[Complete finite-family classification]\label{thm:classification}
For a finite $\Gamma\le S_m$ and a nonempty ground $\Gamma$-set $Y$, the
following are equivalent in $V=M[c]$:
\begin{enumerate}
\item $\CF(\Gamma,Y)$ holds.
\item An admissible rule $D_m(\kappa)\to Y$ belongs to $\ODg$.
\item A nonempty finite family of admissible rules belongs to $\ODg$.
\end{enumerate}
For finite label sets and their action coded by natural numbers, the same
equivalence holds with $\ODlow$ in place of $\ODg$. In the positive case
a rule already belongs to $\OD_{V_\kappa}^V$, without using $q$.
\end{theorem}

\begin{proof}
The implication (2)$\Rightarrow$(3) uses the singleton family. Theorem
\ref{thm:necessity} proves (3)$\Rightarrow$(1). For (1)$\Rightarrow$(2),
fix in $M$ a well-order $\prec$ of the reals and a well-order $\triangleleft$
of $Y$. All finite-alphabet words can be coded by reals in a fixed way.
There are no new reals, so $\prec$ orders all the words under consideration
in $V$ as well.

For each $\Gamma$-orbit $\mathcal O$ of $\sim$-classes of valid words, let
$v_{\mathcal O}$ be the $\prec$-least word whose class belongs to that
orbit. Let $H_{\mathcal O}$ be the stabilizer of $[v_{\mathcal O}]_\sim$.
It is cyclic by Lemma~\ref{lem:stabilizer}. A generator has a fixed point
by $\CF(\Gamma,Y)$; that point is fixed by the whole stabilizer. In the
trivial-stabilizer case every label works. Let $y_{\mathcal O}$ be the
$\triangleleft$-least label fixed by $H_{\mathcal O}$.

For $[w]_\sim=g\cdot[v_{\mathcal O}]_\sim$, set
\[
 F([w]_\sim)=g\cdot y_{\mathcal O}.
\]
If $h$ is another such choice, then $h^{-1}g\in H_{\mathcal O}$, so the
two values agree. This defines an equivariant function of the shift-tail
class. Composing it with $\vec a\mapsto[w_{\vec a}]_\sim$ produces an
admissible rule. It is definable from $\prec,\triangleleft,Y$, and the
action, all ground parameters; it does not use $c$ or $q$.

When $Y$ and the action are finite codes, these parameters apart from
$\prec$ have finite rank, and a relation well-ordering the reals has rank
below $\omega+10$. Thus they all belong to $V_\kappa$, giving the stronger
positive assertion and the corresponding equivalence for $\ODlow$.
\end{proof}

\begin{remark}[Elementwise, not simultaneous, fixed points]
The criterion cannot be strengthened to ``$\Gamma$ has a common fixed
label.'' For example let $\Gamma=S_3$ and take $Y$ to be the disjoint union
of its natural three-point set and its two-point sign-action set. A
transposition fixes a point in the first component; a $3$-cycle fixes both
points in the second. Thus every group element fixes something, but no
label is fixed by the whole group. The positive construction still works.
This example and the underlying cyclic-stabilizer explanation already
appear in the archive. They are useful safeguards against asking more of
the necessity proof than can be true.
\end{remark}

\section{Profinite phases of ultrafilters}\label{sec:phases}

The obstruction for ultrafilters admits a structural strengthening that
is not merely another finite-cycle contradiction. The target is the
profinite completion of the integers,
\[
 \zhat=\varprojlim_{n\mid m}\ZZ/n\ZZ
 =\{(x_n)_{n\ge1}: x_m\bmod n=x_n\text{ whenever }n\mid m\}.
\]
We use its compact Hausdorff topology inherited from the product of the
finite discrete spaces $\ZZ/n\ZZ$.

\subsection{The index cocycle and its residues}

For $a,b\in q$, define an integer
\[
 \chi_a(b)=|b\setminus a|-|a\setminus b|.
\]
For any $a,a',b\in q$, finite cancellation gives
\begin{equation}\label{eq:cocycle}
 \chi_{a'}(b)=\chi_a(b)-\chi_a(a').
\end{equation}
One way to verify this is to take a common finite set outside which all
three representatives agree. The displayed quantities then become
differences of the three finite cardinalities inside that set.
If $a'=a\setminus\{\xi\}$, then $\chi_a(a')=-1$, so
\begin{equation}\label{eq:phaseplus}
 \chi_{a'}(b)=\chi_a(b)+1.
\end{equation}
For $n\ge1$ and $r\in\ZZ/n\ZZ$, define the finite partition
\[
 E^a_{n,r}=\{b\in q:\chi_a(b)\equiv r\pmod n\}.
\]

Here an ultrafilter on $q$ means an ultrafilter on the \emph{full ambient
Boolean algebra} $\powerset(q)^V$, principal or nonprincipal. It is not
the internal tail measure on $\kappa$ discussed later.

Every ultrafilter $W$ on $q$ contains exactly one cell of the $n$th
partition. Call its residue $r_n^a(W)$. The choices are compatible under
divisibility, since the cell at a finer modulus is contained in its
reduction at a coarser one. Hence
\[
 t_a(W)=(r_n^a(W))_{n\ge1}\in\zhat.
\]
Equation~\eqref{eq:phaseplus} implies
\begin{equation}\label{eq:phase-shift}
 t_{a'}(W)=t_a(W)+1\qquad\text{if }a'=a\setminus\{\xi\}.
\end{equation}

\begin{theorem}[Profinite-density theorem]\label{thm:phases}
Let $\mathcal W\in\ODg$ be a nonempty family of ultrafilters on $q$.
Its size is not assumed finite. The phase image
\[
 H_a=\{t_a(W):W\in\mathcal W\}\subseteq\zhat,
 \qquad a=\ran(c),
\]
satisfies $H_a+1=H_a$. Consequently:
\begin{enumerate}
\item $H_a$ is dense in $\zhat$, and its projection onto every
$\ZZ/n\ZZ$ is surjective.
\item $\mathcal W$ cannot be finite.
\item If $\mathcal W$ is compact in the Stone topology on ultrafilters
on $q$, then $H_a=\zhat$ and $|\mathcal W|\ge2^{2^{\aleph_0}}$.
\end{enumerate}
\end{theorem}

\begin{proof}
The space $\zhat$ and its subsets are real-coded low-rank sets. By
Lemma~\ref{lem:lowrank}, $H_a$ is a ground object; more precisely it is an
element of the ground set $(\powerset(\zhat))^M$. Its definition uses
$c$ through $a$, and the fixed definition of $\mathcal W$ from ground
parameters, $q$, and ordinals.

For a late deletion $d$ of $c$, put $a'=\ran(d)$. The family
$\mathcal W$ is unchanged in the same extension. By
\eqref{eq:phase-shift}, $H_{a'}=H_a+1$. Applying
Lemma~\ref{lem:observation} to the phase-image observation gives
$H_{a'}=H_a$, hence $H_a+1=H_a$.

Choose $z\in H_a$. The equality implies $z+\ZZ\subseteq H_a$. This orbit
is dense in $\zhat$: a basic nonempty open set specifies finitely many
compatible congruences, which can be combined into one congruence modulo
the least common multiple of their moduli. An ordinary integer translate
of $z$ satisfies that congruence. In particular each residue modulo $n$
occurs. Since this is true for every $n$, the image and hence
$\mathcal W$ cannot be finite.

For the compactness assertion, the Stone topology has basic clopen sets
$\{W:B\in W\}$ for $B\subseteq q$. Each coordinate map
$W\mapsto r_n^a(W)$ is continuous because its fibres are precisely such
clopen sets, with $B=E^a_{n,r}$. Thus $t_a$ is continuous. The image of
compact $\mathcal W$ in the Hausdorff space $\zhat$ is compact and closed.
Being also dense, it is all of $\zhat$.

Finally $|\zhat|=2^{\aleph_0}$. A compatible residue sequence is a real
code, giving the upper bound. For the lower bound, arbitrary binary digits
determine a $2$-adic component; assign zero residues at every odd prime
power, and use the Chinese remainder theorem to obtain a compatible
sequence at all moduli. This injects $2^\omega$ into $\zhat$.
This explains the first cardinal bound obtainable from the phase map.
For the sharper bound on the family itself, apply
Lemma~\ref{lem:stonecard} below: every infinite compact family of
ultrafilters on a full power set has at least $2^{2^{\aleph_0}}$ members.
The family here is infinite by (2).
\end{proof}

\begin{corollary}[Ultrafilter-valued kernels]\label{cor:ufkernels}
There is no nonempty finite $\ODg$ family of functions that assign to
every $x\in Q_\kappa$ an ultrafilter on $\powerset(x)^V$.
\end{corollary}

\begin{proof}
Evaluate every function in the family at $q$. The resulting set is a
nonempty finite family of ultrafilters on $q$, definable from the same
parameters and $q$, contrary to Theorem~\ref{thm:phases}.
\end{proof}

\subsection{Sharpness: a countable local family exists}

The density conclusion is compatible with countability, and in this case
that compatibility has an explicit realization. It is important to
distinguish a family of ultrafilters on one class from a family of global
kernels on all classes.

\Needspace{12\baselineskip}
\begin{theorem}[The exact local threshold is countable]\label{thm:countable}
Fix a nonprincipal ultrafilter $D$ on $\omega$ in $M$. There is a countably
infinite family $\mathcal U_D(q)$ of nonprincipal ultrafilters on
$\powerset(q)^V$ definable from $D,q$. Its phase image relative to
$a=\ran(c)$ is exactly
\[
 \{t_a(W):W\in\mathcal U_D(q)\}=-\rho(D)+\ZZ\subseteq\zhat,
\]
where $\rho(D)$ records the residues modulo $n$ chosen by $D$ on $\omega$.
The phase map is injective on this family.

Consequently, the least possible cardinality of a nonempty
$\ODlow$ family of ultrafilters on $q$ is exactly $\aleph_0$.
The displayed family is not compact.
\end{theorem}

\begin{proof}
There are no new subsets of $\omega$, so the ground ultrafilter $D$
is an ultrafilter on the full $\powerset(\omega)^V$. For $b\in q$, write
$b^{(k)}$ for the representative obtained by deleting the first $k$
elements of $b$, and define the pushforward ultrafilter
\[
 W_D^b=\{B\subseteq q:\{k<\omega:b^{(k)}\in B\}\in D\}.
\]
Preimages under $k\mapsto b^{(k)}$ preserve complements and finite
intersections, so this is an ultrafilter. It is nonprincipal because the
representatives $b^{(k)}$ are distinct and the preimage of a singleton has
at most one element. Put
\[
 \mathcal U_D(q)=\{W_D^b:b\in q\}.
\]
This set is definable from $q,D$ without choosing a representative of $q$.
We now prove that, despite its apparently $\kappa$-sized indexing set,
it has exactly countably many distinct members.

Fix any $a\in q$ for the proof. Suppose $b,b'\in q$ satisfy
$\chi_a(b)=\chi_a(b')$. Choose a bound above the finite symmetric
difference of $b$ and $b'$. Beyond that bound they have a common tail.
The numbers of their elements before that tail are equal: their difference
is $\chi_b(b')=\chi_a(b')-\chi_a(b)=0$. Thus
\[
 b^{(k)}=(b')^{(k)}\quad\text{for all sufficiently large }k.
\]
For every $B\subseteq q$ the two preimages used to define $W_D^b$ and
$W_D^{b'}$ differ finitely, and $D$ ignores finite differences. Therefore
$W_D^b=W_D^{b'}$. There is at most one ultrafilter in the family for each
integer value of $\chi_a(b)$.

For the converse and the phase computation, let $\rho_n(D)$ be the unique
residue $s\pmod n$ for which $\{k:k\equiv s\pmod n\}\in D$. These
residues are compatible, defining $\rho(D)\in\zhat$. The cocycle identity
gives
\[
 \chi_a(b^{(k)})=\chi_a(b)-k.
\]
It follows directly from the pushforward definition that
\[
 r_n^a(W_D^b)=\chi_a(b)-\rho_n(D)\pmod n,
 \qquad
 t_a(W_D^b)=\chi_a(b)-\rho(D).
\]
Distinct integer indices give distinct elements of $\zhat$: a nonzero
integer is nonzero modulo some positive integer. Thus they give distinct
ultrafilters. Every integer index occurs, because one can delete any
finite number of points from $a$ or add any finite number of new points
to it. There are enough new points since $\kappa$ is uncountable and $a$
is countable. Hence the family is in bijection with $\ZZ$, and its phase
map has exactly the claimed image.

The parameter $D$ belongs to $V_\kappa$. The construction therefore gives
a nonempty countable family in $\ODlow$, while
Theorem~\ref{thm:phases} excludes every nonempty finite family in that
class. This proves the exact threshold. Its countable phase image is
proper and dense in the uncountable Hausdorff space $\zhat$, so it is not
closed. Continuity of the phase map rules out compactness of the family.
\end{proof}

\begin{remark}[Local families versus global kernels]\label{rem:countable}
The construction defines a single assignment
$x\mapsto\mathcal U_D(x)$ on $Q_\kappa$ whose values are countable sets
of ultrafilters. It does not provide a countable definable family of
single-valued kernels $x\mapsto W_x$. Selecting or aligning a labeling of
the local families for all $x$ is an additional problem. Thus the local
countable question is answered positively, but the archive's question
about a countable family of global ultrafilter-valued kernels remains
unsettled here. In particular a definable countable local family need not
have a definable enumeration: such an enumeration would yield a definable
first member, contradicting Theorem~\ref{thm:phases}.
\end{remark}

\subsection{The compact threshold is also attained}

The stronger compactness bound uses the full ambient Boolean algebra
$\powerset(q)$. The following elementary fact is independent of Prikry
forcing. Its proof also fixes the exact cardinal needed in the conclusion.

\begin{lemma}[Infinite compact Stone families]\label{lem:stonecard}
For any set $X$, every infinite compact subspace $F$ of the space of
ultrafilters on $\powerset(X)$ has cardinality at least
$2^{2^{\aleph_0}}$. The space of ultrafilters on a countably infinite set
has cardinality exactly $2^{2^{\aleph_0}}$.
\end{lemma}

\begin{proof}
Write $[A]=\{W:A\in W\}$ for the basic clopen set associated with
$A\subseteq X$. We first construct pairwise disjoint sets $A_n\subseteq X$
and points $W_n\in F\cap[A_n]$. Begin with $B_0=X$ and an infinite trace
$F\cap[B_0]$. Given $B_n$ with infinite trace, separate two distinct
ultrafilters in that trace by some subset of $X$. This partitions $B_n$
into two pieces whose traces in $F$ are both nonempty. At least one trace
is infinite. Retain that piece as $B_{n+1}$ and call the other piece
$A_n$. Choose $W_n$ in its nonempty trace. The $A_n$ are pairwise
disjoint by construction.

For any ultrafilter $D'$ on $\omega$, define
\[
 W_{D'}=\{A\subseteq X:\{n:A\in W_n\}\in D'\}.
\]
It is an ultrafilter because this coordinatewise test preserves Boolean
operations. It lies in the closure of $\{W_n:n<\omega\}$: if
$A\in W_{D'}$, then some $W_n$ contains $A$, and finite intersections
of such basic neighborhoods reduce to a single basic neighborhood.
The compact subspace $F$ is closed in the Hausdorff Stone space, so
$W_{D'}\in F$.

For $S\subseteq\omega$, put $C_S=\bigcup_{n\in S}A_n$.
Disjointness gives $C_S\in W_n$ exactly when $n\in S$. Therefore
\[
 C_S\in W_{D'}\quad\Longleftrightarrow\quad S\in D'.
\]
Distinct ultrafilters $D'$ produce distinct points of $F$. Thus $F$
contains at least as many points as there are ultrafilters on $\omega$.

Here is a proof of the latter cardinality. An ultrafilter is a subset
of $\powerset(\omega)$, so the upper bound is $2^{2^{\aleph_0}}$.
For the lower bound, use the countable set
\[
 Z=\{(n,H):n<\omega,\ H\subseteq 2^n\}.
\]
For each binary sequence $x\in2^\omega$, let
$A_x=\{(n,H):x\restriction n\in H\}$. Given finitely many distinct
$x_1,\ldots,x_k$ and prescribed choices of $A_{x_i}$ or their
complements, their corresponding prefixes are distinct for all
sufficiently large $n$. At every such $n$ one can choose $H$ to meet
exactly the prescribed membership requirements. Thus every such finite
Boolean combination is infinite. For each of the
$2^{2^{\aleph_0}}$ assignments of a zero or one to all the indices
$x\in2^\omega$, the chosen sides, together with the cofinite filter on
$Z$, have the finite intersection property and extend to an ultrafilter.
Different assignments require different ultrafilters. Since $Z$ is
countably infinite, this proves the lower bound on $\omega$ as well.
\end{proof}

\begin{corollary}[The exact compact local threshold]\label{cor:compactsharp}
The least cardinality of a nonempty compact $\ODlow$ family of
ultrafilters on $q$ is exactly $2^{2^{\aleph_0}}$. A family attaining
this bound is the Stone closure
\[
 \mathcal C_D(q)=\overline{\mathcal U_D(q)}.
\]
All its members are nonprincipal, and its phase image is all of $\zhat$.
\end{corollary}

\begin{proof}
The Stone space of ultrafilters on $\powerset(q)$ is compact Hausdorff:
it is the closed space of Boolean homomorphisms to $2$ inside
$2^{\powerset(q)}$. Thus the displayed closure is compact. It is
nonempty and definable from $D,q$; explicitly, membership in the closure
says that every basic neighborhood meets $\mathcal U_D(q)$.
It is infinite since it contains that countably infinite family, so
Lemma~\ref{lem:stonecard} gives the lower bound.

For the matching upper bound, fix $a\in q$ for the proof and let
$R_a=\{a^{(k)}:k<\omega\}$. Every $b\in q$ has a tail ray that
merges with $R_a$, so $R_a\in W_D^b$. Hence the whole countable family
lies in the clopen subspace $[R_a]$, and so does its closure.
Restriction to $R_a$ identifies $[R_a]$ with the ultrafilters on the
countably infinite set $R_a$. This space has cardinality
$2^{2^{\aleph_0}}$ by Lemma~\ref{lem:stonecard}, proving the upper bound.

Every principal ultrafilter is an isolated point, and the set of all
principal ultrafilters is open. Therefore the nonprincipal ultrafilters
form a closed subspace. The family in Theorem~\ref{thm:countable} lies
there, and so does its closure. The phase image of the closure is all
of $\zhat$ by Theorem~\ref{thm:phases}. Finally any nonempty compact
definable family is infinite by that same theorem, and hence has at least
the displayed cardinality by Lemma~\ref{lem:stonecard}. This proves
minimality among all such definable families, not only for the example.
\end{proof}

The positive constructions in Theorem~\ref{thm:countable} and
Corollary~\ref{cor:compactsharp} do not intrinsically require Prikry
forcing: they work for every class in $Q_\lambda$, for any uncountable
$\lambda$ of cofinality $\omega$, using a nonprincipal ultrafilter on the
ambient $\powerset(\omega)$. The Prikry hypothesis supplies the
nonexistence of smaller definable families at the generic class.

\section{Charges: the exact finite-additivity contrast}\label{sec:charges}

A \emph{charge} on $q$ is a finitely additive probability measure
$\mu:\powerset(q)^V\to[0,1]$. Countable additivity is not included in the
definition. We first describe a restriction on definable finite families
and then construct examples showing why the restriction is not stronger.

\begin{theorem}[Residue laws for finite families of charges]\label{thm:charges}
Suppose $\mathcal M\in\ODg$ is a nonempty family of charges on $q$ and
$h=|\mathcal M|<\omega$. Let $a=\ran(c)$.
\begin{enumerate}
\item If $h=1$, its member $\mu$ satisfies
$\mu(E^a_{n,r})=1/n$ for all $n\ge1$ and every residue $r$.
\item For general $h$, every $\mu\in\mathcal M$ satisfies
$\mu(E^a_{\ell,r})=1/\ell$ for every prime $\ell>h$ and every residue $r$.
\item Every member has infinitely many values, gives measure zero to each
integer-index fibre $\chi_a^{-1}\{k\}$, and fails countable additivity.
\end{enumerate}
\end{theorem}

\begin{proof}
For each modulus $n$, form the set of probability vectors
\[
 H_n(a)=\{(\mu(E^a_{n,r}))_{r<n}:\mu\in\mathcal M\}\subseteq[0,1]^n.
\]
It is nonempty, has size at most $h$, and belongs to $M$: each entry is a
real, there are no new reals, and the set of vectors is finite. On deleting
one point of $a$, the $r$th coordinate becomes the old $(r-1)$st
coordinate. The deletion lemma makes $H_n(a)$ invariant under this cyclic
rotation.

When $h=1$, the one vector is fixed by the rotation, so all its coordinates
are equal. Their sum is $1$ by finite additivity, hence they are all $1/n$.
For a prime modulus $\ell>h$, an orbit of a vector under rotation has
length either $1$ or $\ell$. Since the invariant set of vectors has at
most $h<\ell$ elements, every vector is fixed. Again its coordinates must
all be $1/\ell$.

For a fixed $\mu\in\mathcal M$, the values $1/\ell$ for primes $\ell>h$
are infinitely many distinct members of its range. For each integer $k$,
\[
 \chi_a^{-1}\{k\}\subseteq E^a_{\ell,k\bmod\ell},
 \qquad
 0\le\mu(\chi_a^{-1}\{k\})\le1/\ell.
\]
Letting these primes grow shows that this fibre has measure zero. The
integer-index fibres form a countable partition of $q$. Countable
additivity would therefore give $\mu(q)=0$, contradicting $\mu(q)=1$.
\end{proof}

In particular no nonempty finite definable family can consist of
finite-range charges, or of countably additive probability measures.
Ultrafilters correspond to the two-valued charges and are an immediate
special case, although Theorem~\ref{thm:phases} says substantially more
about arbitrary ultrafilter families.

\subsection{An explicit positive kernel}

The supplied synthesis already proves the existence of a definable
charge-valued kernel from a shift-invariant mean. The following is a
concrete Ces\`aro--ultrafilter implementation of that positive construction,
not a new claim of existence.

Fix a nonprincipal ultrafilter $D$ on $\omega$ in $M$. Such a $D$ exists
by Choice, lies in $V_\kappa$, and remains an ultrafilter on all subsets of
$\omega$ in $V$ because no new reals are added.
For $a\in D_\kappa$, let $a^{(k)}$ be $a$ with its first $k$ elements
deleted. If $x=[a]$ and $B\subseteq x$, define
\begin{equation}\label{eq:positivekernel}
 K_D(x)(B)=\lim_{N\to D}\frac1N
       \bigl|\{k<N:a^{(k)}\in B\}\bigr|.
\end{equation}
The ultrafilter limit of a bounded real sequence is the unique limit along
$D$ in the compact interval $[0,1]$.

\begin{proposition}[A definable finitely additive kernel]\label{prop:positive}
Formula~\eqref{eq:positivekernel} is independent of the chosen
representative $a\in x$. It defines a charge on every $x\in Q_\kappa$,
and the entire kernel $K_D$ is ordinal definable from $D$. In particular,
finitely additive probability kernels exist in $\OD_{V_\kappa}^V$.
\end{proposition}

\begin{proof}
If $a=^*b$, choose $r,s<\omega$ such that
$a^{(r)}=b^{(s)}$. Then $a^{(r+j)}=b^{(s+j)}$ for every $j<\omega$.
For each $B\subseteq[a]$, the two sequences of indicator values along
these rays have a common tail after shifting their indices by a fixed
integer. Their first $N$ averages differ by at most a constant divided by
$N$. This tends to zero along the nonprincipal ultrafilter $D$, so the
limits agree.

For disjoint $B,C$, the indicator of their union is the sum of the two
indicators at every point of the ray. Taking finite averages and then the
ultrafilter limit preserves this equality. The empty set has value $0$ and
the whole class has value $1$, while all values are nonnegative. Thus the
formula defines a charge. Independence permits the definition to quantify
over a representative without making a choice of one; the result is a
single function of $x$ and $D$. Replacement gives the kernel as a set,
and its definition uses only $D$ and the ordinal $\kappa$.
\end{proof}

This proof uses no embedding hypothesis. The ray-average method works for
any uncountable $\lambda$ of cofinality $\omega$ in a universe where the
chosen $D$ is an ultrafilter on its full $\powerset(\omega)$. At the generic
class it has the uniform residue values required by
Theorem~\ref{thm:charges}. It also assigns zero to each singleton: the ray
$a^{(0)},a^{(1)},\ldots$ has distinct terms, so a singleton can occur at
most once in an average.

\subsection{Sharpness for finite families}

It would be incorrect to infer from Theorem~\ref{thm:charges} that every
member of a finite definable family is uniform at \emph{every} modulus.
There are exact finite examples showing the difference between a singleton
and an unordered family.

\begin{proposition}[Periodic ray averages]\label{prop:periodiccharges}
For every integer $d\ge1$ there is a $d$-element family of charges on $q$
definable from $q,D$. When $a=\ran(c)$ is used to label its members as
$\mu_0^a,\ldots,\mu_{d-1}^a$, one has
\[
 \mu_r^a(E^a_{d,-r})=1.
\]
For a prime $d$, this shows that the threshold $\ell>h$ in
Theorem~\ref{thm:charges}(2) cannot be replaced by $\ell\ge h$.
\end{proposition}

\begin{proof}
For any $a\in q$ put
\[
 \mu_r^a(B)=\lim_{N\to D}\frac1N
       \bigl|\{j<N:a^{(dj+r)}\in B\}\bigr|,
 \qquad 0\le r<d.
\]
Each is a charge by the same finite-additivity argument as above.
If $a^{(r_0)}=b^{(s_0)}$, then their full tail rays agree after the fixed
shift $s_0-r_0$. The subray with indices congruent to $r$ modulo $d$
therefore agrees, after a finite shift of its own index, with the subray
of $b$ having residue $r+s_0-r_0$ modulo $d$. Finite shifts do not affect
the averages. Thus changing $a$ permutes the set
$\{\mu_r^a:r<d\}$, and that unordered set is definable from $q,D$.

Since $\chi_a(a^{(dj+r)})=-(dj+r)$, every term in this subray belongs
to $E^a_{d,-r}$. The corresponding charge assigns this cell value $1$
and every other cell of the $d$th partition value $0$. The $d$ charges
are therefore distinct. For prime $d\ge2$, none is uniform at modulus
$d$, even though the definable family has exactly $d$ members.
\end{proof}

The last construction is a finite family of \emph{local charges on one
class}. It does not choose a labeling of these charges simultaneously
across all classes and therefore should not be silently reinterpreted as
a finite family of global kernels. The singleton kernel in
Proposition~\ref{prop:positive} has no such ambiguity.

\section{The ground trace, rigidity, and the canonical core}\label{sec:core}

The measurable-strength rigidity and normal-trace constructions are already
in the supplied synthesis \cite[Section 13]{archive}. Here we prove the
specific version needed to put the new finite-choice conclusions in the
\emph{same} model. This section also shows that those conclusions do not
force the normal trace to concentrate on smaller measurables.

\subsection{Homogeneity relative to the generic class}

\begin{lemma}[Ground containment for definable sets of ordinals]\label{lem:groundOD}
Every set of ordinals in $\ODg$ belongs to $M$.
\end{lemma}

\begin{proof}
First consider a sentence with ground parameters, ordinal parameters, and
the canonical name for the generic finite-difference class. Any two
conditions can be strengthened to conditions whose cones are isomorphic
by replacing the stems.

Explicitly, given $(s,A)$ and $(t,B)$, shrink the upper parts to a common
$C\in U$ contained in $A\cap B$ and above both stems. Below $(s,C)$,
conditions have stems $s$ followed by a finite sequence $u$ from $C$.
Replace that stem by $t$ followed by the same $u$, leaving the later upper
part unchanged. This is an isomorphism with the cone below $(t,C)$.
Corresponding generic presentations have the form $s$ followed by a common
infinite tail and $t$ followed by that tail. They generate the same
extension and have the same finite-difference class. Ground parameters
are interpreted identically. Consequently the original sentence has the
same truth value in the corresponding presentations.

This is the usual forcing-invariance argument for isomorphic cones,
with invariance of the class parameter checked explicitly. In particular,
two conditions cannot force opposite truth values for such a sentence.
Since deciding conditions are dense, the maximum condition decides every
such sentence.

Now let $x\subseteq\eta$ be uniquely definable in $V$ from allowed
parameters. The assertion that there is a unique such subset of $\eta$
is true in $V$ and is therefore forced by the maximum condition, by the
preceding paragraph. For each $\alpha<\eta$, the maximum condition also
decides whether $\alpha$ belongs to that uniquely defined set. Thus
\[
 x=\{\alpha<\eta:\mathbf1_{\PP_U}\Vdash
          \alpha\text{ belongs to the uniquely defined subset}\}.
\]
The forcing relation for this fixed formula and its ground parameters is
definable in $M$. Separation in $M$ therefore puts $x$ in $M$.
Any set of ordinals is bounded by some ordinal, so the claim follows.
\end{proof}

\begin{lemma}[Hereditary ground containment]\label{lem:HODground}
Both
\[
 N=\HOD_{V_\kappa\cup\{q\}}^V
 \quad\text{and}\quad H_p=\HOD_{\{p,q\}}^V\quad(p\in M)
\]
are subclasses of $M$.
\end{lemma}

\begin{proof}
Use induction on rank in either hereditary class. Suppose $x$ belongs to
that class and all its elements belong to $M$ by induction. Choose an
ordinal $\beta$ with $x\subseteq V_\beta^M$, and choose in $M$ an
ordinal-indexed enumeration $e:\theta\to V_\beta^M$. Such an enumeration
exists by Choice in $M$. The set
\[
 I_x=\{\xi<\theta:e(\xi)\in x\}
\]
is a set of ordinals definable from $q$ and ground parameters. Here the
parameters formerly allowed from $V_\kappa^V$ belong to $M$ by
Lemma~\ref{lem:lowrank}; the additional parameter $e$ is also in $M$.
Lemma~\ref{lem:groundOD} puts $I_x$ in $M$. Hence $x=e[I_x]\in M$.
The empty-rank case starts the induction.
\end{proof}

\subsection{Regularity and rigidity}

The class $N$ above is the hereditary definability inner model used in the
synthesis; it satisfies $\ZF$ and contains $V_\kappa$. The latter
inclusion follows directly because every element of $V_\kappa$, and every
member of its transitive closure, is an allowed parameter. Choice is not
needed in $N$ for the following regularity argument.

\begin{proposition}[A rigid generic class]\label{prop:rigidity}
The cardinal $\kappa$ is regular in $N$. Equivalently in the terminology
of the synthesis, $q$ is rigid: no nonempty family of fewer than $\kappa$
cofinal subsets of $\kappa$ of order type below $\kappa$ belongs to
$\ODlow$. Also, for $T\in\ODlow$ with $T\subseteq D_\kappa$,
\[
 |T\cap q|\in\{0,\kappa\}.
\]
\end{proposition}

\begin{proof}
A cofinal map from an ordinal below $\kappa$ to $\kappa$ belonging to $N$
would belong to $M$ by Lemma~\ref{lem:HODground}. It would still be
cofinal there, contradicting regularity of the measurable cardinal
$\kappa$ in $M$. Thus $\kappa$ is regular in $N$.

Suppose $\mathcal A\in\ODlow$ is a nonempty family of the indicated
kind with $|\mathcal A|<\kappa$. Let $\mu$ be the least order type occurring
among its members, and let $\mathcal A_\mu$ consist of the members of that
order type. Their union $u$ is cofinal in $\kappa$ and is definable from
the allowed parameters. Moreover
\[
 |u|\le |\mathcal A_\mu|\cdot |\mu|<\kappa.
\]
This inequality uses only a product of two fixed cardinals below the
infinite cardinal $\kappa$, not an invalid regularity assumption about
$\kappa$ in $V$. Because $u\subseteq\kappa$, its order type is below
$\kappa$. As a definable set of ordinals, $u$ belongs to $N$, and its
increasing enumeration there contradicts regularity. Conversely a short
cofinal subset in $N$ would give a forbidden singleton family; this also
explains the equivalence with the archive's rigidity formulation.

Every class $q$ has cardinality $\kappa$. Fix $a\in q$. Every member is
obtained from $a$ by finitely many additions and deletions, giving the
upper bound $\kappa$. Adding one point $\alpha\in\kappa\setminus a$
gives $\kappa$ distinct members, giving the lower bound. If $T\cap q$
were nonempty of size below $\kappa$, it would be a forbidden definable
family of cofinal $\omega$-sets. This proves the last statement.
\end{proof}

\subsection{The internal normal measure}

Choose in $M$ a bijection $b:\kappa\to V_\kappa^M$, using the
inaccessibility of $\kappa$ there, and set
\[
 H=\HOD_{\{b,q\}}^V,
 \qquad
 U_H=\{A\in\powerset(\kappa)\cap H:\tail(q,A)\}.
\]
Unlike the many-parameter model $N$, the hereditary model with the two
fixed set parameters $b,q$ satisfies $\ZFC$: the usual definition-code
well-order of ordinal-definable objects gives Choice in this hereditary
class. We use that standard interpretation of $\HOD_{\{b,q\}}$ throughout.

\begin{theorem}[A normal tail measure in the canonical core]\label{thm:core}
The following hold:
\begin{enumerate}
\item $H\subseteq M$, $V_\kappa^H=V_\kappa^V$, and $b\in H$.
\item $U_H\in H$, and $H$ regards it as a normal nonprincipal
$\kappa$-complete ultrafilter on $\kappa$. Externally, $U_H=U\cap H$.
\item Every representative of $q$ is Prikry generic over $H$ for $U_H$;
all give the same extension $H[c]$, with
\[
 H\subsetneq H[c]\subseteq V,
 \qquad V_\kappa^{H[c]}=V_\kappa^V.
\]
The extension $H[c]$ has the same cardinals as $H$, contains $q$, and has
cofinality $\omega$ at $\kappa$. No representative of $q$ belongs to $H$.
\end{enumerate}
\end{theorem}

\begin{proof}
Lemma~\ref{lem:HODground} gives $H\subseteq M$. For every $x\in
V_\kappa^M$, choose $\alpha<\kappa$ with $x=b(\alpha)$. This makes $x$
definable from $b$ and an ordinal. The same is true of every member of its
transitive closure, because all those members also lie in $V_\kappa^M$.
Thus $x\in H$. Together with Lemma~\ref{lem:lowrank} this gives equality
of the rank segments. The graph $b$ is itself definable from $b$, and
its elements and their descendants consist of ordinals and these same
low-rank sets. Hence $b\in H$.

For $A\in\powerset(\kappa)\cap H$, we have $A\in M$. If $A\in U$,
Mathias' criterion gives $a_c\subseteq^*A$. If $A\notin U$, its complement
belongs to $U$, so $a_c$ is almost disjoint from $A$ and cannot be almost
contained in $A$. Therefore
\[
 \tail(q,A)\quad\Longleftrightarrow\quad A\in U,
\]
which proves $U_H=U\cap H$. The definition of $U_H$ from $b,q$ makes it
ordinal definable from those parameters, and every one of its members is
in $H$. Therefore it is hereditarily so definable, and $U_H\in H$.
It is an ultrafilter on the subsets of $\kappa$ in $H$, and is
nonprincipal, by the corresponding properties of $U$.

For internal completeness, take a sequence $\langle A_i:i<\delta\rangle
\in H$, where $\delta<\kappa$ and all $A_i\in U_H$. The entire sequence
belongs to $M$. Completeness of $U$ in $M$ puts its intersection in $U$;
that intersection belongs to $H$, so it belongs to $U_H$. For normality,
take in $H$ a regressive function on a set in $U_H$. This is a ground
regressive function on a member of $U$. Normality in $M$ provides a
constant fibre in $U$. The fibre is defined from the function and its
ordinal value, so it belongs to $H$, hence to $U_H$. These are precisely
the internal completeness and normality assertions.

Every $a\in q$ is cofinal of order type $\omega$ and is almost contained
in every $A\in U_H$, by definition. Mathias' criterion, now over $H$,
proves its genericity. Finite changes yield the same extension because
the finitely many extra ordinals belong to $H$. Prikry preservation over
$H$ gives equality of cardinals and the new cofinality. Every finite
modification of $a$ can be formed in $H[a]$; conversely every member of
$q$ is such a modification. Therefore its internally computed class is
exactly the ambient $q$, so $q\in H[a]$. Since $H$ already contains the
full ambient $V_\kappa$, that rank segment stays equal in the intermediate
model. Finally $H$ sees $\kappa$ measurable and hence regular; a
representative of $q$ would give it a cofinal $\omega$-sequence. Thus no
representative belongs to $H$, and the inclusion $H\subsetneq H[c]$
is proper.
\end{proof}

\begin{assessment}{Two different ultrafilters, with different domains}
The internal normal measure $U_H$ is an ultrafilter on
$\powerset(\kappa)\cap H$, and its completeness is tested on sequences
belonging to $H$. The forbidden ultrafilters of
Theorem~\ref{thm:phases} live on $\powerset(q)^V$. These statements do not
conflict. In fact the ambient countable sequence of tail intervals of
$c$ demonstrates that internal $\kappa$-completeness of $U_H$ is not
ambient countable completeness.
\end{assessment}

\begin{corollary}[The new obstructions do not force concentration]\label{cor:firstmeas}
Suppose $\kappa$ is the least measurable cardinal of $M$. Then $H$ regards
$\kappa$ as its least measurable cardinal, while all the finite-family
conclusions above hold in $V$. In particular $U_H$ does not concentrate
on the cardinals below $\kappa$ that are measurable in $H$.
\end{corollary}

\begin{proof}
For $\alpha<\kappa$, the objects needed to express that $\alpha$ carries
a nonprincipal $\alpha$-complete ultrafilter all lie in $V_\kappa$:
subsets of $\alpha$, candidate ultrafilters on them, and sequences of
length below $\alpha$ of their elements have rank bounded below the
inaccessible ordinal $\kappa$. The models $H$ and $M$ agree on this entire
rank segment. They therefore agree about measurability below $\kappa$.
There are no such measurables in $M$, so there are none in $H$.
Theorem~\ref{thm:core} makes $\kappa$ measurable in $H$. The set of
internally measurable cardinals below it is empty and cannot be in
$U_H$. All earlier arguments required only that $U$ be a normal measure
on $\kappa$, so they still apply.
\end{proof}

This is a separation \emph{at the displayed witness}: the new finite
choice failures do not recover the concentration-on-measurables feature
of the ultraexacting critical-sequence construction in the archive. It
is not a claim that no stronger normal-trace witness exists anywhere else
in the universe.

\section{The precise consistency calibration}\label{sec:consistency}

We give a statement without the external ground-model symbol $M$.
This prevents a conditional model construction from being confused with
an unqualified new lower-bound theorem for bare selector failure.

\begin{definition}[The finite cyclic obstruction at a class]\label{def:FC}
For a strong limit uncountable cardinal $\lambda$ of cofinality $\omega$
and $q\in Q_\lambda$, let $\Inv(\lambda,q)$ assert the following.
For every finite $\Gamma\le S_m$ and every nonempty finite $\Gamma$-set
$Y$ coded by natural numbers, if $\CF(\Gamma,Y)$ fails, then there is no
nonempty finite family of admissible rules $D_m(\lambda)\to Y$ in
$\OD_{V_\lambda\cup\{q\}}$.
\end{definition}

The sufficiency half of the finite-label classification needs no extra
axiom here: the incidence-word construction uses a well-order of the
reals, which is an available parameter in $V_\lambda$. Thus
$\Inv(\lambda,q)$ includes the full equivalence between the cyclic
condition, existence of a definable rule, and existence of a nonempty
finite definable family. It also implies every finite-selector
obstruction in Corollary~\ref{cor:selectors}.

\begin{definition}[The calibrated package]\label{def:package}
Let $\NTFC$ assert that there exist $\lambda,b,q$ such that:
\begin{enumerate}
\item $\lambda$ is an uncountable strong limit cardinal with
$\cf(\lambda)=\omega$, $q\in Q_\lambda$, and $b:\lambda\to V_\lambda$
is a bijection;
\item $\lambda$ is regular in $\HOD_{V_\lambda\cup\{q\}}$;
\item in $H=\HOD_{\{b,q\}}$, the set
$\{A\in\powerset(\lambda)\cap H:\tail(q,A)\}$ belongs to $H$ and is a
normal nonprincipal $\lambda$-complete ultrafilter on $\lambda$;
\item $\Inv(\lambda,q)$ holds.
\end{enumerate}
These are assertions about definable inner classes in the usual language
of set theory. Relative ordinal definability and membership in the
hereditary classes have their standard uniform first-order definitions;
no satisfaction predicate for the whole universe is being added.
\end{definition}

\begin{theorem}[Exact measurable-strength package]\label{thm:consistency}
Over $\ZFC$,
\[
 \Con(\ZFC+\text{there is a measurable cardinal})
 \quad\Longleftrightarrow\quad
 \Con(\ZFC+\NTFC).
\]
The equivalence remains true if $\NTFC$ also requires that the displayed
$\lambda$ be the \emph{least} measurable cardinal of its displayed core
$\HOD_{\{b,q\}}$.
\end{theorem}

\begin{proof}
For the upper bound, start with a measurable cardinal and a normal measure
in a ground model of $\ZFC$. Force with its ordinary Prikry forcing.
Lemma~\ref{lem:lowrank} supplies the strong limit singular cardinal and
the ground bijection $b:\kappa\to V_\kappa$. Proposition
\ref{prop:rigidity} supplies regularity in the many-parameter hereditary
model. Theorem~\ref{thm:core} supplies the normal tail measure in the
fixed-parameter core. Finally Theorem~\ref{thm:necessity}, using
$\ODlow\subseteq\ODg$, supplies $\Inv(\kappa,q)$. Therefore the extension
satisfies $\NTFC$. Choosing the least measurable cardinal in the ground
and applying Corollary~\ref{cor:firstmeas} gives the additional clause.

For the lower bound, any model of the package has the definable inner
model $H=\HOD_{\{b,q\}}$ satisfying $\ZFC$. Clause (3) says exactly that
$H$ has a measurable cardinal, witnessed by the displayed internal
ultrafilter. Thus it yields a model of the measurable-cardinal theory.

These are the standard relative-consistency arguments by set forcing and
by interpretation in a definable inner model. The transitive-model
notation used to present the mathematics is not an extra hypothesis on
the two consistency statements. Equivalently one can apply the forcing
interpretation and the inner-model interpretation to each finite proof
that might witness inconsistency.
\end{proof}

\begin{remark}[Where the lower bound comes from]\label{rem:lowerbound}
The lower bound in Theorem~\ref{thm:consistency} comes from the normal-trace
clause, not from $\Inv(\lambda,q)$ alone. The archive already calibrated
normal traces at measurable strength. The contribution is that adding all
the formerly ultraexacting finite cyclic obstructions does not increase
that strength. No theorem here identifies the exact consistency strength
of the bare failure of definable binary choice on a cofinal quotient.
\end{remark}

The upper-bound model also has the profinite-density theorem and the
finite-family charge restrictions from Sections~\ref{sec:phases} and
\ref{sec:charges}. Those are consequences of the forcing construction;
they need not be included in the definition of the package to obtain its
stated calibration.

\section{Consequences for the research program}\label{sec:conclusions}

\textbf{A resolved forcing question.}
The Prikry-extension part of the synthesis's last question group (ii) is
settled by Theorems~\ref{thm:necessity} and \ref{thm:classification} and
Corollary~\ref{cor:selectors}. The result is uniform in every ordinary
normal-measure Prikry extension, not just in a specially selected model.
The allowed definition parameters include the invariant generic class
$q$ and arbitrary ground sets. This is stronger than merely removing
ultraexactingness from the $\OD_{V_\kappa}$ singleton conclusion.

\textbf{A distinction between two parts of the ultraexacting picture.}
Finite cyclic failures and the existence of one normal-trace core can
coexist when that core has no smaller measurable cardinal. They therefore
do not reproduce, at the same witness, the concentration-on-measurables
property established for the ultraexacting critical-sequence trace in the
archive. The latter remains a genuinely additional property in this
comparison. Nothing here transfers the archive's internal successor
measurability theorem to the present Prikry construction.

\textbf{A structural obstruction rather than only a cardinal bound.}
For ultrafilters, the conclusion is the exact dynamical relation
$H_a+1=H_a$ on a canonically extracted phase image. Density, all finite
residue surjections, and the compactness lower bound follow from it.
Theorem~\ref{thm:countable} realizes a single dense integer orbit and
shows that the least size of a nonempty definable local family is exactly
$\aleph_0$. Corollary~\ref{cor:compactsharp} gives the exact compact
threshold $2^{2^{\aleph_0}}$, also with a definable example.
For charges, finite families really can exist, and the periodic ray
examples exhibit the sharp obstruction at prime moduli. Thus replacing
ultrafilters by finitely additive measures changes the answer in a way
that is both constructive and quantitatively visible.

\begin{question}[Rigidity without a Prikry presentation]
Does regularity of $\lambda$ in $\HOD_{V_\lambda\cup\{q\}}$, by itself,
imply $\Inv(\lambda,q)$? The proof here needs a late-deletion principle
for ground-valued observations. No derivation of that principle from
regularity alone is supplied.
\end{question}

\begin{question}[Countable global families]
Can a countably infinite $\OD_{V_\kappa}$ family of single-valued
ultrafilter-valued kernels on $Q_\kappa$ exist? The countable local
families of Theorem~\ref{thm:countable} do not furnish a countable family
of global sections. The analogous question for countable families of
selectors also remains outside the present proof.
\end{question}

\begin{question}[A wider forcing criterion]
Which other singularizing forcings admit both a same-extension deletion
operation and enough ground-valued response tables to reproduce
Theorem~\ref{thm:necessity}? An application must verify those properties
for the actual forcing; no blanket assertion is made here for all
Prikry-type or tree-Prikry forcings.
\end{question}

\appendix
\section{Finite checks and a proof audit}\label{app:checks}

The archive accompanying this report contains \nolinkurl{finite_check.py},
a dependency-free Python program, and its actual output in
\nolinkurl{finite_check_results.txt}. The program checks four elementary
parts of the construction: the root identity in
Lemma~\ref{lem:root}; the sign and index in the deletion formula using
finite integer sequences with disjoint blocks; the absence of a fixed
proper subset under a full cycle; and the $S_3$ example with elementwise
but no simultaneous fixed points. It also checks the finite residue
frequencies behind the periodic ray averages.

These tests are useful for catching sign errors in $g^{-1}$ or
$\tau^{-n}$. They are not tests of a generic extension, of the forcing
theorem, of hereditary ordinal definability, or of any large-cardinal
consistency assertion. No Lean formalization of the new arguments is
included, and none is claimed.

\begin{center}
\small
\renewcommand{\arraystretch}{1.15}
\begin{tabularx}{\textwidth}{@{}P{.33\textwidth}Y@{}}
\toprule
\textbf{Potential gap}&\textbf{How it is avoided}\\
\midrule
Replacing a definable finite family by a definable member&The proof uses the unordered set of all finite response tables.\\
Pretending that deletion is an automorphism of $V$&The forcing truth lemma compares two generic presentations of the same universe.\\
Using a new high-rank response as a ground constant&Only finite tables with ground entries, or genuinely unchanged low-rank objects, are decided.\\
Confusing the generic class with a representative&The class $q$ is invariant; $c$ is used only in the observation and in test configurations.\\
Using regularity at the ambient singular $\kappa$&Small unions are bounded using two fixed cardinal factors; regularity is used only in the inner model.\\
Equating internal and ambient measures&$U_H$ is on internal subsets of $\kappa$; forbidden ultrafilters are on full ambient subsets of $q$.\\
Overstating the consistency result&Its lower bound is explicitly supplied by the normal-trace clause.\\
Confusing the two size thresholds&The local minimum is $\aleph_0$; the compact minimum is $2^{2^{\aleph_0}}$, using a separate Stone-space argument.\\
\bottomrule
\end{tabularx}
\end{center}

\section{Source and notation concordance}\label{app:concordance}

The notation $D_\lambda,Q_\lambda,\tail,\CF$ follows the supplied
\emph{Synthesis}. Its ``rigidity is regularity'' theorem is used here in
the explicit small-family form of Proposition~\ref{prop:rigidity}.
Its fixed-parameter normal-trace model $\HOD_{\{b,q\}}$ is denoted by
$H$ here; $M$ is reserved for the actual Prikry ground. Our $V=M[c]$ is
the ambient extension, not the internal Prikry core $H[c]$.

The synthesis's finite-label sufficiency proof is expanded in
Section~\ref{sec:sufficiency}. Its measurable-strength trace construction
is expanded in Section~\ref{sec:core}. Neither reproduction is offered
as a new theorem. The new forcing transfer is
Theorem~\ref{thm:necessity}; the main further refinements are
Theorems~\ref{thm:phases}, \ref{thm:countable}, and \ref{thm:charges},
Corollary~\ref{cor:compactsharp}, Proposition~\ref{prop:periodiccharges},
and the compatibility/separation
statements in Corollary~\ref{cor:firstmeas} and
Theorem~\ref{thm:consistency}.

The supplied Lean development includes files such as
\nolinkurl{Cardinals/FiniteLabel.lean},
\nolinkurl{Cardinals/Combinatorics/FiniteCycles.lean}, and
\nolinkurl{Cardinals/Combinatorics/ThirdRound.lean}. Their role in this
continuation is terminological and conceptual reference. The finite
permutation tests in this archive are a separate, limited check and do
not change the verification status of that development.

\begin{thebibliography}{9}
\small

\bibitem{archive}
\emph{Large cardinals at the inconsistency frontier: a synthesis}.
User-supplied research archive \nolinkurl{Cardinals4.zip},
\nolinkurl{docs/Large_Cardinals_Synthesis.tex}, dated 18 September 2026;
with \nolinkurl{docs/Large_Cardinals_Unified_Report.tex} and the
\nolinkurl{Cardinals/} Lean reference development. Especially the sections
``The exact finite-label classification,'' ``Measures hidden in rigid
classes,'' ``Exact consistency strength,'' and the final open-question
group for reports 19--27. Unpublished supplied material.

\bibitem{benhamou}
Tom Benhamou.
\emph{Prikry forcing and tree Prikry forcing of various filters}.
Archive for Mathematical Logic \textbf{58} (2019), 787--817.
\href{https://doi.org/10.1007/s00153-019-00660-3}{doi:10.1007/s00153-019-00660-3}.
Preprint: \href{https://arxiv.org/abs/1801.04424}{arXiv:1801.04424};
theorem numbers in this report refer to version 2, 25 September 2018.

\bibitem{mathias}
A.~R.~D. Mathias.
\emph{On sequences generic in the sense of Prikry}.
Journal of the Australian Mathematical Society \textbf{15} (1973),
409--414. The normal-measure criterion is also restated and proved as
Theorem 3.17 in \cite{benhamou}.

\bibitem{abl}
Juan P. Aguilera, Joan Bagaria, and Philipp L\"ucke.
\emph{Large cardinals, structural reflection, and the HOD Conjecture}.
\href{https://arxiv.org/abs/2411.11568}{arXiv:2411.11568}.
Background for the exacting/ultraexacting program; not an input to the
new deletion proof.

\bibitem{abgl}
Juan P. Aguilera, Joan Bagaria, Gabriel Goldberg, and Philipp L\"ucke.
\emph{Large cardinals beyond HOD}.
\href{https://arxiv.org/abs/2509.10254}{arXiv:2509.10254}, version 1,
12 September 2025.
Background connecting large embeddings, HOD, and Prikry constructions.

\bibitem{goldberg}
Gabriel Goldberg.
\emph{Consistency beyond the Kunen inconsistency}.
Lecture notes reporting joint work with Doug Blue,
Warsaw Inner Model Theory Meeting, 1 July 2026,
especially Remark 2.4.
\href{https://math.berkeley.edu/~goldberg/Slides/CoverExact.pdf}{Author-hosted PDF}.

\end{thebibliography}
\end{document}
